We conduct a methodology-locked cross-market replication of the volatility-based
momentum--reversal switching rule proposed for equities by \citet{butt2023}. The setting
is a broad cross-section of Chinese commodity and financial futures. The mapping uses
causal main-contract selection, same-contract returns, an 11-month momentum signal that
skips the most recent month, a one-month reversal signal, and a five-year rolling
volatility distribution. The strategy holds momentum unless lagged market volatility is
in its highest quartile, in which case it holds reversal. After the state-history
requirement, 66 monthly observations are evaluable. Under a lagged open-interest-notional
market proxy, the switch has a higher endpoint than momentum, but its mean monthly
increment is economically small and statistically indistinguishable from zero. Reversal
beats momentum in only 4 of 14 high-volatility months, and one event determines the sign
of the average advantage. Under an equal-product market proxy, high-volatility momentum
outperforms reversal and the switch underperforms momentum. The volatility proxies are
correlated, yet their quartile assignments agree in only 42.4\% of holding months.
High-volatility months exhibit higher cross-product correlation, dispersion, and
directional coherence, showing that volatility captures shock intensity. It does not,
however, uniquely identify shock direction, source, breadth, or stage. The evidence
therefore rejects a transportable volatility-only router in this sample without rejecting
momentum, reversal, or state-dependent allocation more generally.
